%% ma: L12-B18-0016
%% lop: 12
%% chuong: 6
%% bai: 18
%% dang: 12.18.7
%% loai: TN4
%% muc: TH
%% dap_an: "C"
%% nguon: "(NBV)-10 ĐỀ KIỂM TRA HK2-TOÁN 12-DUNG CHUNG 3 SÁCH.pdf (H02, Đề số 2), Phần 1 Câu 11"
%% kien_thuc: "Bài 18 (công thức nhân, xác suất của biến cố đối có điều kiện)"
%% trang_thai: chinh-thuc
\begin{ex}
Cho các biến cố $A$ và $B$ sao cho $0<P(A)<1$, $0<P(B)<1$. Khẳng định nào sau đây là sai?
\choice{$P(B)=P(AB)+P(\overline AB)$}{$P(A\mid B)=1-P(\overline A\mid B)$}{\True $P(\overline AB)+P(A\overline B)=1$}{$P(AB)=P(A)\cdot P(B\mid A)$}
\loigiai{A đúng vì $B=AB\cup\overline AB$ (hai biến cố xung khắc). B đúng vì $P(A\mid B)+P(\overline A\mid B)=1$. D đúng theo công thức nhân. C sai: $P(\overline AB)+P(A\overline B)$ là xác suất xảy ra đúng một trong hai biến cố, nói chung nhỏ hơn $1$ (chẳng hạn khi $A=B$ thì bằng $0$). Chọn C.}
\end{ex}
